\chapter{Experiment files}\label{app:files}

The AEDL files of the C GLoBES distribution are included in
\file{test/data}, together with their flux, cross section and migration matrix
files. They describe the experiments as they were planned when the files were
written and should be checked against current designs before use.

\begin{center}\small
\begin{tabularx}{\linewidth}{lX}
\toprule
File & Experiment \\
\midrule
\file{NOvA.glb} & NOvA: $\numu$ beam from NuMI, 25~kt detector at 812~km, $\nue$
  appearance and $\numu$ disappearance in neutrino and antineutrino mode \\
\file{NOvA\_pre\_post\_smearing.glb} & the same, with efficiencies as pre- and
  post-smearing lists \\
\file{T2K.glb} & T2K: J-PARC beam to Super-Kamiokande (22.5~kt) at 295~km \\
\file{T2HK.glb} & T2HK: upgraded beam to a 500~kt water Cherenkov detector at 295~km \\
\file{NFstandard.glb} & neutrino factory: 50~GeV muons, 50~kt magnetised iron
  detector at 3000~km \\
\file{NFvar.glb} & the same with the AEDL variables \code{BASELINE} and \code{emax} \\
\file{NF\_hR\_lT.glb} & neutrino factory with high energy resolution and low
  threshold (variables \code{BASELINE}, \code{emax}) \\
\file{BB\_350.glb} & beta beam with $\gamma=350$, 500~kt water Cherenkov
  detector at 730~km, with migration matrices from a detector simulation \\
\file{BBvar\_WC.glb} & beta beam with variable $\gamma$ (\code{gammafactor},
  \code{baselinefactor}, \code{EXP\_FACTOR}) \\
\file{Reactor1.glb}, \file{Reactor2.glb} & reactor experiments with 20 and
  100~t detectors at 1.7~km \\
\file{D-Chooz\_near.glb}, \file{D-Chooz\_far.glb} & Double Chooz near (0.1~km)
  and far (1.05~km) detectors, for a near--far analysis with a user-defined
  $\chi^2$ (\file{examples/example5.jl}) \\
\file{Simple.glb} & the example of Chapter~\ref{ch:aedl-intro} \\
\bottomrule
\end{tabularx}
\end{center}

\noindent Cross sections: \file{XCC.dat} (charged current), \file{XNC.dat}
(neutral current), \file{XQE.dat} (quasi-elastic), \file{XCCreactor.dat}
(inverse beta decay). Fluxes: \file{NOvAplus.dat}, \file{NOvAminus.dat},
\file{JHFplus.dat}, \file{JHFminus.dat}, \file{Reactor.dat},
\file{BB350flux.dat}.

\chapter{Examples}\label{app:examples}

The directory \file{examples} contains
\begin{description}[style=nextline,leftmargin=3.6cm]
\item[\file{tour.jl}] the tour of Chapter~\ref{ch:tour}, with the high-level API;
\item[\file{example-tour.jl}] the same with the \code{glb*} functions, a line-by-line
  translation of \file{example-tour.c};
\item[\file{example1.jl}] the $\sst{13}$--$\dcp$ correlation with systematics only;
\item[\file{example2.jl}] projections onto the $\sst{13}$ axis;
\item[\file{example3.jl}] locating the sign degeneracy and its shape;
\item[\file{example4.jl}] the $\theta_{13}$ sensitivity with statistics,
  systematics, correlations and degeneracies;
\item[\file{example5.jl}] a reactor experiment with near and far detectors
  (Double Chooz) and user-defined $\chi^2$ functions;
\item[\file{example6.jl}] a user-defined probability engine: wave packet
  decoherence in a reactor experiment.
\end{description}
\file{examples/README.md} compares their output with the C examples.

\chapter{Building this manual}

The manual is in \file{docs/manual}. Its figures are drawn by pgfplots from
tables computed by GLoBES.jl. To rebuild everything:
\begin{julia}
cd docs/manual
make            # runs figures/make_figures.jl and examples/tour.jl, then lualatex
\end{julia}
\code{make figures} recomputes only the data, \code{make pdf} only typesets.
